\documentclass[12pt,letterpaper]{article}
\usepackage[utf8]{inputenc}
\usepackage[spanish]{babel}
\spanishdeactivate{<>}
\usepackage{amsmath,amssymb}
\usepackage{geometry}
\usepackage{booktabs}
\usepackage{xcolor}
\usepackage[hidelinks]{hyperref}
\usepackage{tikz}
\usetikzlibrary{babel}

\geometry{top=2.0cm, bottom=2.0cm, left=2.0cm, right=2.0cm}

\title{\textbf{Análisis Matemático de Cinemática Directa e Inversa}\\
\large Brazo Robótico $T^3$ Community Arm (3 Eslabones Planares)}
\author{Documentación de Tesis ROS 2}
\date{\today}

\begin{document}
\maketitle

\begin{center}
\begin{tikzpicture}[scale=0.33]
    % Ejes principales
    \draw[thick] (-3,0) -- (31,0) node[right] {$R$};
    \draw[thick] (0,0) -- (0,28) node[above] {$Z$};
    
    % Base del robot
    \draw[fill=gray!30] (-1.5,0) rectangle (1.5,13);
    \node at (0,6.5) {$L_{\text{base}}$};
    
    % Eje vertical inferior de referencia desde el hombro (q2 = 0 deg)
    \draw[dashed, blue, line width=1.4pt] (0,13) -- (0,1) node[below] {Vertical Inferior ($q_2=0^\circ$)};
    
    % Hombro
    \filldraw[fill=orange!80] (0,13) circle (0.8) node[left=10pt] {\textbf{Hombro}};
    
    % Eslabón 1 (L1)
    \draw[line width=4pt, orange!90!black] (0,13) -- (9.9,22.9) node[midway, above left] {$L_1$};
    
    % Arco q2 (medido desde el eje vertical INFERIOR hacia arriba/derecha)
    \draw[line width=1.5pt, blue] (0,8) arc (-90:45:5);
    \node[blue] at (5.5, 9.5) {\Large \textbf{$q_2$}};
    
    % Codo
    \filldraw[fill=orange!80] (9.9,22.9) circle (0.8) node[above=8pt] {\textbf{Codo}};
    
    % Prolongación recta de L1 (referencia para q3 = 0 deg)
    \draw[dashed, gray, line width=1.2pt] (9.9,22.9) -- (16.96,29.96) node[above right] {Prolongación de $L_1$};
    
    % Eslabón 2 (L2)
    \draw[line width=4pt, orange!90!black] (9.9,22.9) -- (23.05,18.11) node[midway, above right] {$L_2$};
    
    % Arco q3 (medido desde la prolongación de L1)
    \draw[line width=1.5pt, red] (14.14,27.14) arc (45:-20:6);
    \node[red] at (17.5, 25) {\Large \textbf{$q_3$}};
    
    % Muñeca
    \filldraw[fill=black] (23.05,18.11) circle (0.6) node[above] {\textbf{Muñeca}};
    
    % Offsets de pinza extending FORWARD to the RIGHT (+R)
    \draw[line width=2pt, darkgray] (23.05,18.11) -- (28.52,18.11) node[midway, above] {$d_x$};
    \draw[line width=2pt, darkgray] (28.52,18.11) -- (28.52,15.94) node[midway, right] {$d_z$};
    
    % Efector Final
    \filldraw[fill=red] (28.52,15.94) circle (0.7) node[below right] {\textbf{Efector Final} $(r, Z)$};
\end{tikzpicture}\\
\small \textbf{Figura 1:} Esquema geométrico del manipulador en el plano $RZ$. El ángulo del hombro $q_2$ se mide desde la \textbf{Vertical Inferior ($-Z$)}. El ángulo del codo $q_3$ se mide relativo a la prolongación del eslabón $L_1$. La pinza extiende hacia la derecha ($+d_x$).
\end{center}

\section{Definición Geométrica y Parámetros del Robot}

El robot $T^3$ Community Arm se modela en el plano de elevación $RZ$ como un manipulador antropomórfico planar de 3 eslabones principales con base rotacional ($q_1$) y dos eslabones de elevación ($q_2, q_3$).

\begin{table}[h!]
\centering
\begin{tabular}{@{}lll@{}}
\toprule
\textbf{Parámetro} & \textbf{Símbolo} & \textbf{Valor} \\ \midrule
Altura de la Base & $L_{\text{base}}$ & $13.0\text{ cm}$ \\
Longitud Eslabón 1 (Hombro) & $L_1$ & $14.0\text{ cm}$ \\
Longitud Eslabón 2 (Codo) & $L_2$ & $14.0\text{ cm}$ \\
Offset de Pinza Longitudinal & $d_x$ & $+5.47\text{ cm}$ \\
Offset de Pinza Vertical & $d_z$ & $-2.17\text{ cm}$ \\ \bottomrule
\end{tabular}
\caption{Parámetros dimensionales del robot Community Arm.}
\end{table}

\subsection{\texorpdfstring{Definición de Ángulos Articulares ($q_2, q_3$)}{Definicion de Angulos Articulares (q2, q3)}}

\begin{itemize}
    \item $q_1$: Rotación de la base en el plano $XY$ (Azimut).
    \item $q_2$: Ángulo de la articulación del hombro, medido desde la \textbf{Vertical Inferior ($-Z$)}.
    \item $q_3$: Ángulo de la articulación del codo, medido relativo a la \textbf{prolongación del eslabón $L_1$}.
\end{itemize}

\section{Cinemática Directa (Forward Kinematics - FK)}

Dado el vector de ángulos articulares $q = (q_1, q_2, q_3)$ expresado en radianes:

\subsection{Conversión a Ángulos Absolutos respecto a la Horizontal}

Para proyectar los eslabones trigonométricamente sobre el sistema cartesiano $+R$ (horizontal), se definen los ángulos absolutos $\theta_2$ y $\theta_3$:

\begin{equation}
\theta_2 = q_2 - \frac{\pi}{2}
\end{equation}

\begin{equation}
\theta_3 = \theta_2 - q_3 = q_2 - q_3 - \frac{\pi}{2}
\end{equation}

\subsection{Ecuaciones de Posición de la Muñeca}

En el plano de elevación $RZ$:
\begin{equation}
r_{\text{wrist}} = L_1 \cos(\theta_2) + L_2 \cos(\theta_3)
\end{equation}
\begin{equation}
z_{\text{wrist}} = L_{\text{base}} + L_1 \sin(\theta_2) + L_2 \sin(\theta_3)
\end{equation}

\subsection{Ecuaciones Finales del Efector Final (Gripper Tip)}

Debido al paralelogramo horizontal rígido, los offsets $(d_x, d_z)$ actúan como traslaciones puras:
\begin{equation}
r_{\text{total}} = r_{\text{wrist}} + d_x
\end{equation}
\begin{equation}
X = r_{\text{total}} \cdot \cos(q_1) = \left( L_1 \cos(\theta_2) + L_2 \cos(\theta_3) + d_x \right) \cos(q_1)
\end{equation}
\begin{equation}
Y = r_{\text{total}} \cdot \sin(q_1) = \left( L_1 \cos(\theta_2) + L_2 \cos(\theta_3) + d_x \right) \sin(q_1)
\end{equation}
\begin{equation}
Z = L_{\text{base}} + L_1 \sin(\theta_2) + L_2 \sin(\theta_3) + d_z
\end{equation}

\section{Cinemática Inversa (Inverse Kinematics - IK)}

Dado un punto objetivo $(X, Y, Z)$ en coordenadas cartesianas globales, se calcula $q = (q_1, q_2, q_3)$:

\subsection{\texorpdfstring{Paso 1: Ángulo de la Base ($q_1$)}{Paso 1: Angulo de la Base (q1)}}
\begin{equation}
r_{xy} = \sqrt{X^2 + Y^2}
\end{equation}
\begin{equation}
q_1 = \text{atan2}(Y, X)
\end{equation}

\subsection{Paso 2: Aislamiento de la Muñeca}
\begin{equation}
r_{\text{wrist}} = r_{xy} - d_x
\end{equation}
\begin{equation}
z' = Z - d_z - L_{\text{base}}
\end{equation}

\subsection{\texorpdfstring{Paso 3: Ley de Cosenos para el Ángulo del Codo ($\delta$)}{Paso 3: Ley de Cosenos para el Angulo del Codo (delta)}}
\begin{equation}
D^2 = r_{\text{wrist}}^2 + (z')^2
\end{equation}
\begin{equation}
\cos(\delta) = \frac{r_{\text{wrist}}^2 + (z')^2 - L_1^2 - L_2^2}{2 L_1 L_2}
\end{equation}
\begin{equation}
\delta = \text{atan2}\left(\pm \sqrt{1 - \cos^2(\delta)}, \; \cos(\delta)\right)
\end{equation}

\subsection{\texorpdfstring{Paso 4: Solución de $q_2$ y $q_3$}{Paso 4: Solucion de q2 y q3}}
\begin{equation}
A = L_1 + L_2 \cos(\delta), \quad B = L_2 \sin(\delta)
\end{equation}
\begin{equation}
\theta_2 = \text{atan2}\left( A \cdot z' - B \cdot r_{\text{wrist}}, \; A \cdot r_{\text{wrist}} + B \cdot z' \right)
\end{equation}
\begin{equation}
q_2 = \theta_2 + \frac{\pi}{2}
\end{equation}
\begin{equation}
q_3 = -\delta
\end{equation}

\end{document}
